\section{Introduction}
%
%\GD{Following is just to help me structure my thoughts}
%\begin{itemize}
%\item 3DGS and variants use a mixture of splats and SH's to represent appearance
%\item Texture is a compact and effective way of representing high frequency
%appearance attached to a surface
%\item Total number of parameters of the representation: Primitives, with position (mean), scale, rotation for shape together with opacity and SH for appearance; and texture resolution.
%\item We strive to exploit compact texture as much as possible; this requires a choice of how to distribute the parameter burden between the primitives and the texture in a principled manner.
%\item We introduce a principled approach to resolution and primitive management, based on the content in the input images, by setting world-space texel size based on available resolution in the input.
%\end{itemize}
%
%\TODO{Standard Intro}

Gaussian Splatting~\cite{3dgs,2dgs} has become the method of choice for the  capture and real-time novel-view synthesis of real scenes. It relies on flexible, easy-to-optimize \emph{Gaussian primitives} to represent the scene.
However, in the standard approach, fine visual appearance -- such as intricate texture on a surface -- needs to be represented with a large number of very small primitives even when the geometric complexity is low. 
This wasteful limitation arises because each Gaussian primitive is associated with a single color sample, preventing the decoupling of appearance and geometric complexity.
Rendering techniques, on the other hand, have traditionally built on \emph{texture mapping} to represent such detailed appearance. We build on this tradition to propose a representation for {\em textured} Gaussian primitives that is guided by the scene and its content.

In recent and concurrent work~\cite{bbsplat, gstex, textured3dgs, supergaussians} a number of methods introduce textured Gaussian primitives, 
providing a spatially varying color.
Most of these methods
use planar 2D Gaussians to match the dimensionality of traditional textures.
Each primitive carries several parameters (position, rotation, scale, opacity and spherical harmonics (SH)),
while textures are 2D arrays of RGB(A) texels.
The fundamental challenge of such approaches is how to allocate model capacity across the scene
and between the number of texels and the number of primitives used.
Previous and concurrent solutions either deal with the issue by imposing a fixed texture resolution per primitive~\cite{bbsplat, supergaussians, textured3dgs}, or propose heuristic approaches, that are far from ideal~\cite{gstex}.
This leaves open the question of how to share resources in a more principled manner.%between texels and primitives.

To answer this question,
we propose a \emph{content-aware} texturing solution for 2D Gaussian primitives that uses the appearance complexity \revision{rev6315 term content unclear}{as observed from the input images of the scene} to make informed decisions on how parameter resources
are distributed.
We do this by first introducing a textured Gaussian primitive representation in which texel size is \emph{fixed} in world space.
\revision{rev6315 term content unclear}{
This makes the appearance texture of each primitive independent of its size,
and therefore unaffected by the growing or shrinking that it undergoes during optimization.
This representation separates the parameters of geometry and appearance, allowing us to refine them independently with appropriate allocation of memory capacity.
%As a result, we can develop methods that allocate model capacity to refine the two independently.
%This enables us to refine 
}
% , estimated using the content of the input images.
More specifically,
we introduce a method to increase and decrease the texel size as optimization evolves,
by analyzing the frequency content that primitive textures can represent and determining the corresponding error.
We pair that with a resolution-aware method to control the overall number of primitives,
striking a balance between texture size and number of primitives.
Our solution interacts closely with the optimization,
using the downscale/upscale mechanism to reduce the error in \emph{appearance},
while our control of the number of primitives handles \emph{geometric} error by spawning additional Gaussians to capture geometric detail.

%\begin{itemize}
%\item 3DGS is great, but requires tons of small primitives, mixing appearance and shape 
%\item Texture is great (standard CG), good representation of surface detail
%\item Several efforts to combine the two. Problem of how to assign texture resolution, usually ad-hoc
%\item We introduce a principled, content-aware approach
%\begin{itemize}
%\item Review previous methods at a high level, say heuristics
%\item Need for a principled approach rather than ad-hoc
%\end{itemize}
%
%\begin{itemize}
%\item We introduce a principled content-aware approach, 
%\item Explain representation, need for frequency content adaptation and
%progressive approach, need for splitting
%\item Contribs: 

\noindent
In summary,
we propose three contributions:
\begin{itemize}
\item A texture representation for 2D Gaussians that defines texels with a \emph{fixed} world space size, supporting visual reconstruction independent of the shape of the primitives.
% \item A content-aware texture representation for Gaussian primitives,
\item A progressive algorithm that adaptively determines texel size and allows for content-aware fitting of the scene.
\item A resolution-based solution to control the number of primitives,
adjusted to the textured representation.

\end{itemize}
Our experiments show that our method provides a good balance between visual quality and the number of parameters used for the representation, comparing favorably to other textured Gaussian primitives solutions proposed in recent and concurrent work.